\begin{afterword}
\summaryhead

The post takes the anti-zombie argument of the previous posts as settled: consciousness is, among other things, the cause of our talk about consciousness, so it cannot be removed while everything physical stays the same. It then asks how far the argument generalizes. A switch said to end your consciousness without affecting your brain is ruled out, because you would go on talking about consciousness for the same reasons.

Since every physical change has some effect (the switch pulls on your atoms by gravity), the post concedes that the general principle is about expectations, not certainty. But an effect far smaller than thermal noise should not be expected to switch off consciousness, and neither should a word heard across the room. The principle: a change that cannot alter your internal narrative should not make you cease to be conscious.

The post ends with a dialogue about replacing your neurons one by one with functionally identical robots. Albert says you would still be you; Bernice says you would be dead; Charles says a new conscious person would replace you. The author sides with Albert and defers the reasons.

\responsehead

The post is careful about the limits of its argument. It sets the zombie question aside openly. It concedes that its definition of consciousness settles no empirical question and is ``not a knockdown argument,'' and that the switch argument is about expectations, not certainty. Its physics checks: the pull of a one-gram switch ten metres away is about $6.7 \times 10^{-16}$ m/s$^2$, far below thermal noise.

What the post does not give is the history of its central case. Replacing neurons one by one with functionally identical parts is a standard thought experiment, discussed by Pylyshyn, Searle and others before 1995. Searle described Bernice's outcome: experience shrinking to nothing while behaviour stays the same. Chalmers answered with the argument the post gives Albert. A change the system cannot notice is not a change in experience, since ``whenever experiences change significantly and we are paying attention, we can notice the change.'' Chalmers even used a switch, between a neural and a silicon circuit. So the leading advocate of the zombie argument already held Albert's view on consciousness, for Albert's reason. That is support for the post, and a reader should know it.

The principle, as this post states it, does not yet settle the robot case, and the post says as much. It covers changes after which the causes of your talk are ``pretty much the same.'' The robots keep the talk and replace its causes; whether that falls under the principle is Charles's objection. The reasons for Albert's view came in later posts that this book does not include. What the post establishes is the switch and the spoken word.

Two smaller points. The post describes the anti-zombie argument as an argument against epiphenomenalism; Chalmers held that the zombie argument does not require epiphenomenalism, as the notes on ``Zombie Responses'' explain. And the reason given for moving on is that most of the blog's commenters agree, which the post itself does not treat as settling a ``philosophically controversial'' question.

In short: a careful, openly hedged extension of the anti-zombie argument to small physical changes, whose closing case, gradual neuron replacement, is a known thought experiment in which the zombie argument's chief advocate had already argued for the post's side.
\end{afterword}
